\chapter{Inverse and Implicit Function Principles}
\label{ch:ii09}

The inverse and implicit function principles are local nonlinear analogues of solving an invertible linear system. This chapter is a fresh canonical construction because the source atlas contains no mapped legacy chapter for II/09.

\section*{Learning goals}
After this chapter, the reader should be able to:
\begin{itemize}
\item verify the hypotheses of the inverse function theorem by computing the derivative and checking its invertibility at a specified point;
\item distinguish local invertibility from global injectivity and construct examples where the former holds without the latter;
\item compute the derivative of a local inverse from \(D(f^{-1})(f(a))=Df(a)^{-1}\);
\item identify which variables can be solved for in \(F(x,y)=0\) by checking the appropriate derivative block;
\item derive and evaluate the implicit-derivative formula \(Dg=-D_yF^{-1}D_xF\) in concrete systems;
\item analyze failure cases where a singular derivative or rank loss prevents the theorem from applying, without concluding more than the hypotheses justify;
\end{itemize}
\section*{Conceptual roadmap}
\[
\boxed{\text{definition}\;\longrightarrow\;\text{estimate}\;\longrightarrow\;\text{theorem}\;\longrightarrow\;\text{application}.}
\]

\section{Local invertibility}
A nonlinear map with an invertible derivative should behave like its linearization on a sufficiently small neighborhood.

\section{Inverse function theorem}
Invertibility of $Df(a)$ gives a local differentiable inverse near $a$.

\section{Derivative of the inverse}
The inverse derivative is the inverse linear map: $D(f^{-1})(f(a))=Df(a)^{-1}$.


% BEGIN VOL02-EXPANSION II09-example-01
\begin{example}[A genuinely nonlinear local inverse with a computable inverse derivative]\label{ex:ii09-expand-01}
Define
\[
f(x,y)=(x+y^2,\ y+x^2).
\]
At the origin,
\[
Df(0,0)=I,
\]
so the inverse function theorem gives a local \(C^1\) inverse near \(0\).
Moreover
\[
D(f^{-1})(0)=Df(0,0)^{-1}=I.
\]
To see the nonlinear content, note that solving \(f(x,y)=(u,v)\) is not a
linear system: \(x=u-y^2\), \(y=v-x^2\). The theorem guarantees that these
coupled equations have one nearby solution for every nearby \((u,v)\), and
that the solution depends smoothly on the data.
\end{example}
% END VOL02-EXPANSION II09-example-01
\section{Implicit equations}
A system $F(x,y)=0$ can be solved locally for $y$ as a function of $x$ when the $y$-Jacobian is invertible.


% BEGIN VOL02-EXPANSION II09-example-02
\begin{example}[Solving a coupled implicit system for two variables]\label{ex:ii09-expand-02}
Let
\[
F(t,x,y)=
\begin{pmatrix}
x+y-t\\
x-y^2
\end{pmatrix}.
\]
At \((t,x,y)=(0,0,0)\),
\[
D_{(x,y)}F=
\begin{pmatrix}
1&1\\
1&0
\end{pmatrix},
\qquad \det=-1\ne0.
\]
Hence there are unique smooth functions \(x(t),y(t)\) near \(0\) solving
\(F=0\). Differentiating,
\[
D_{(x,y)}F\binom{x'}{y'}+
D_tF=0,\qquad D_tF=\binom{-1}{0}.
\]
Solving at \(0\) gives \(x'(0)=0\), \(y'(0)=1\). This is the matrix-valued
implicit derivative formula in a concrete system.
\end{example}
% END VOL02-EXPANSION II09-example-02
\section{Implicit differentiation}
The derivative of the implicit function is obtained by solving the differentiated linear system.

\section{Regular level sets}
Full-rank derivatives make level sets locally look like Euclidean subspaces.

\section{Contraction proof strategy}
One proof route reformulates the nonlinear equation as a contraction on a small ball.

\section{Failure modes}
Singular derivative, noninjectivity, or rank loss can destroy local solvability or uniqueness.


% BEGIN VOL02-EXPANSION II09-example-03
\begin{example}[A singular derivative means theorem failure, not automatic impossibility]\label{ex:ii09-expand-03}
Consider
\[
F(x,y)=y^3-x.
\]
At the origin \(F_y(0,0)=0\), so the implicit function theorem cannot solve
smoothly for \(y\) as a function of \(x\) there. Nevertheless the equation has
the unique continuous solution
\[
y=x^{1/3}.
\]
The solution is not differentiable at the origin. Thus a singular derivative
means the theorem's smooth-local-solvability conclusion is unavailable; it
does not by itself mean that no solution or no uniqueness exists.
\end{example}
% END VOL02-EXPANSION II09-example-03
\section{Core structural results}

\begin{theorem}[Inverse function theorem]\label{thm:ii09-01}
If $f:\mathbb R^n\to\mathbb R^n$ is $C^1$ near $a$ and $Df(a)$ is invertible, then $f$ is a $C^1$ diffeomorphism between neighborhoods of $a$ and $f(a)$.
\begin{proof}
After composing with $Df(a)^{-1}$, reduce to derivative near the identity. On a sufficiently small ball the nonlinear error has Lipschitz constant below one, so solving $f(x)=y$ becomes a contraction problem. Differentiability of the inverse follows from the linearized equation.
\end{proof}
\end{theorem}

\begin{theorem}[Derivative of inverse]\label{thm:ii09-02}
Under the hypotheses of the inverse theorem, $D(f^{-1})(f(a))=Df(a)^{-1}$.
\begin{proof}
Differentiate $f^{-1}\circ f=\mathrm{id}$ and use the chain rule.
\end{proof}
\end{theorem}

\begin{theorem}[Implicit function theorem]\label{thm:ii09-03}
If $F:\mathbb R^{n+m}\to\mathbb R^m$ is $C^1$, $F(a,b)=0$, and $D_yF(a,b)$ is invertible, then near $(a,b)$ the equation $F(x,y)=0$ has a unique $C^1$ solution $y=g(x)$.
\begin{proof}
Apply the inverse function theorem to $\Phi(x,y)=(x,F(x,y))$, whose derivative is block triangular with invertible diagonal blocks.
\end{proof}
\end{theorem}

\section{Worked examples}

\begin{example}[One-dimensional inverse theorem]\label{ex:ii09-worked-01}
Explain why $f'(a)\ne0$ yields local invertibility for a $C^1$ real function. Continuity keeps $f'$ of one sign near $a$, so $f$ is strictly monotone there. The inverse exists locally and differentiating $f(f^{-1}(y))=y$ gives the inverse derivative.
\end{example}

\begin{example}[A locally invertible map]\label{ex:ii09-worked-02}
For $f(x,y)=(e^x\cos y,e^x\sin y)$, determine where the derivative is invertible. The Jacobian determinant is $e^{2x}$, always positive, so the derivative is invertible everywhere; global injectivity fails because of $2\pi$-periodicity in $y$.
\end{example}

\begin{example}[Local versus global inverse]\label{ex:ii09-worked-03}
Why does the inverse function theorem not imply a global inverse for the previous map? The theorem is local. Distinct points differing by $(0,2\pi k)$ have the same image despite nonsingular derivative everywhere.
\end{example}

\begin{example}[Implicit circle graph]\label{ex:ii09-worked-04}
Near $(0,1)$ solve $x^2+y^2=1$ for $y$ and compute $y'(0)$. Take $F=x^2+y^2-1$; $F_y=2y=2$ at $(0,1)$, so $y=\sqrt{1-x^2}$ locally and $y'=-x/y$, hence $y'(0)=0$.
\end{example}

\section{Solved dossiers}
The dossiers are longer solved problems. Legacy mappings guide topic selection where explicit retained source blocks exist; otherwise fresh canonical problems complete the chapter.

\begin{problem}[One-dimensional inverse theorem]\label{prob:ii09-dossier-01}
Explain why $f'(a)\ne0$ yields local invertibility for a $C^1$ real function.
\end{problem}
\begin{solution}
Continuity keeps $f'$ of one sign near $a$, so $f$ is strictly monotone there. The inverse exists locally and differentiating $f(f^{-1}(y))=y$ gives the inverse derivative.
\end{solution}

\begin{problem}[A locally invertible map]\label{prob:ii09-dossier-02}
For $f(x,y)=(e^x\cos y,e^x\sin y)$, determine where the derivative is invertible.
\end{problem}
\begin{solution}
The Jacobian determinant is $e^{2x}$, always positive, so the derivative is invertible everywhere; global injectivity fails because of $2\pi$-periodicity in $y$.
\end{solution}

\begin{problem}[Local versus global inverse]\label{prob:ii09-dossier-03}
Why does the inverse function theorem not imply a global inverse for the previous map?
\end{problem}
\begin{solution}
The theorem is local. Distinct points differing by $(0,2\pi k)$ have the same image despite nonsingular derivative everywhere.
\end{solution}

\begin{problem}[Implicit circle graph]\label{prob:ii09-dossier-04}
Near $(0,1)$ solve $x^2+y^2=1$ for $y$ and compute $y'(0)$.
\end{problem}
\begin{solution}
Take $F=x^2+y^2-1$; $F_y=2y=2$ at $(0,1)$, so $y=\sqrt{1-x^2}$ locally and $y'=-x/y$, hence $y'(0)=0$.
\end{solution}

\begin{problem}[Implicit derivative formula]\label{prob:ii09-dossier-05}
Derive $Dg(a)=-D_yF(a,b)^{-1}D_xF(a,b)$.
\end{problem}
\begin{solution}
Differentiate $F(x,g(x))=0$: $D_xF+D_yF\,Dg=0$, then solve for $Dg$.
\end{solution}

\begin{problem}[Regular level curve]\label{prob:ii09-dossier-06}
Show $x^2+y^2=1$ is locally a smooth curve at every point.
\end{problem}
\begin{solution}
The gradient $(2x,2y)$ never vanishes on the unit circle, so one partial derivative is nonzero and the implicit theorem applies.
\end{solution}

\begin{problem}[A singular failure]\label{prob:ii09-dossier-07}
What happens to $F(x,y)=y^2-x^3=0$ at the origin?
\end{problem}
\begin{solution}
$D_yF(0,0)=0$, so the theorem does not apply. The cusp cannot be represented as a single smooth function $y=g(x)$ through the origin.
\end{solution}

\begin{problem}[Coordinate straightening]\label{prob:ii09-dossier-08}
Explain geometrically what the implicit theorem says about a regular level set.
\end{problem}
\begin{solution}
After a local $C^1$ change of coordinates, the equation becomes the coordinate condition $y=0$, so the level set is locally flat.
\end{solution}

\begin{problem}[Newton linearization link]\label{prob:ii09-dossier-09}
Relate inverse-theorem solvability to Newton's method.
\end{problem}
\begin{solution}
Both use invertibility of the derivative to solve a linearized correction equation. The theorem supplies local uniqueness and stability; Newton iterates attempt to realize that correction computationally.
\end{solution}

\begin{problem}[Parameter dependence]\label{prob:ii09-dossier-10}
If $F(x,y)=y^3+y-x$, show $y$ is a smooth local function of $x$ everywhere on the solution set.
\end{problem}
\begin{solution}
$F_y=3y^2+1>0$ everywhere, so the implicit theorem applies at every solution point.
\end{solution}

\begin{problem}[Rank perspective]\label{prob:ii09-dossier-11}
Why is the determinant condition coordinate-dependent notation for an invariant statement?
\end{problem}
\begin{solution}
Nonzero determinant means the derivative linear map is an isomorphism. Under basis changes the determinant is multiplied by nonzero factors, so invertibility itself is invariant.
\end{solution}

\begin{problem}[IFT synthesis]\label{prob:ii09-dossier-12}
Summarize inverse and implicit theorems as nonlinear linear algebra.
\end{problem}
\begin{solution}
An invertible derivative lets one solve all variables locally; an invertible derivative block lets one solve selected variables locally. Both are perturbative versions of solving nonsingular linear systems.
\end{solution}

\section{Exercises with complete solutions}

\begin{exercise}\label{exr:ii09-01}
Can $x\mapsto x^3$ use the inverse theorem at $0$?
\end{exercise}
\begin{hint}
Check the derivative. Compute \(Df(a)\) first and check whether it is an isomorphism; in coordinates this is the nonzero-Jacobian-determinant test.
\end{hint}
\begin{solution}
No; $f'(0)=0$, although a continuous inverse exists globally.
\end{solution}

\begin{exercise}\label{exr:ii09-02}
Find the inverse derivative formula in one dimension.
\end{exercise}
\begin{hint}
Differentiate an identity. Differentiate \(f^{-1}\circ f=\mathrm{id}\) at the base point and solve the resulting linear identity for \(D(f^{-1})\).
\end{hint}
\begin{solution}
$(f^{-1})'(f(a))=1/f'(a)$.
\end{solution}

\begin{exercise}\label{exr:ii09-03}
For $F(x,y)=x+y+xy$, compute $F_y(0,0)$.
\end{exercise}
\begin{hint}
Differentiate in $y$. For \(F(x,y)=0\), compute the derivative only with respect to the variables you want to solve for and check that block for invertibility.
\end{hint}
\begin{solution}
It is $1$, so $y$ is locally a function of $x$.
\end{solution}

\begin{exercise}\label{exr:ii09-04}
Solve the previous equation explicitly.
\end{exercise}
\begin{hint}
Factor in $y$. After verifying \(D_yF\) is invertible, differentiate \(F(x,g(x))=0\) and solve the linear equation for \(Dg\).
\end{hint}
\begin{solution}
$y=-x/(1+x)$ near $0$.
\end{solution}

\begin{exercise}\label{exr:ii09-05}
Why does nonsingular Jacobian imply local openness?
\end{exercise}
\begin{hint}
Use local diffeomorphism. A nonsingular derivative gives only local invertibility; test global injectivity separately, for example by looking for periodicity or symmetry.
\end{hint}
\begin{solution}
A local diffeomorphism maps neighborhoods to neighborhoods.
\end{solution}

\begin{exercise}\label{exr:ii09-06}
State a failure of global injectivity with nonsingular derivative.
\end{exercise}
\begin{hint}
Use periodicity. If the derivative is singular, conclude only that the theorem does not apply; then inspect the equation directly before claiming nonexistence or nonuniqueness.
\end{hint}
\begin{solution}
The exponential-polar map is a standard example.
\end{solution}

\begin{exercise}\label{exr:ii09-07}
What derivative block matters for solving $F(x,y)=0$ for $y$?
\end{exercise}
\begin{hint}
Differentiate with respect to the solved variables. For a regular level set, check surjectivity of the full derivative and identify which coordinate block can be used to graph the set locally.
\end{hint}
\begin{solution}
$D_yF$.
\end{solution}

\begin{exercise}\label{exr:ii09-08}
What is a regular value condition?
\end{exercise}
\begin{hint}
Think surjectivity. To connect the theorem with contraction mapping, normalize the derivative to the identity and estimate the nonlinear remainder on a sufficiently small ball.
\end{hint}
\begin{solution}
At every point of the level set, the derivative is surjective.
\end{solution}


% BEGIN VOL02-EXPANSION II09-exercises-01
\section{Graded supplementary exercises}
These exercises extend the preserved foundational set and are graded by mathematical role.

\subsection*{Standard computations}

\begin{exercise}[Jacobian invertibility]\label{exr:ii09-expand-01}
For \(f(x,y)=(x+y,x-y)\), compute \(Df\) and decide whether the inverse theorem applies.
\end{exercise}
\begin{hint}
Compute the determinant.
\end{hint}
\begin{solution}
\(Df=\begin{pmatrix}1&1\\1&-1\end{pmatrix}\) with determinant \(-2\), so it applies everywhere.
\end{solution}

\begin{exercise}[Inverse derivative]\label{exr:ii09-expand-02}
If \(Df(a)=\begin{pmatrix}2&1\\0&3\end{pmatrix}\), compute \(D(f^{-1})(f(a))\).
\end{exercise}
\begin{hint}
Invert the triangular matrix.
\end{hint}
\begin{solution}
\(\begin{pmatrix}1/2&-1/6\\0&1/3\end{pmatrix}\).
\end{solution}

\begin{exercise}[Implicit slope]\label{exr:ii09-expand-03}
For \(F(x,y)=x^2+xy+y^2-3\), compute \(dy/dx\) where \(F_y\ne0\).
\end{exercise}
\begin{hint}
Use \(y'=-F_x/F_y\).
\end{hint}
\begin{solution}
\(y'=-(2x+y)/(x+2y)\).
\end{solution}

\begin{exercise}[Solve-for block]\label{exr:ii09-expand-04}
For \(F:\mathbb R^3\to\mathbb R\), which derivative must be nonzero to solve \(F(x,y,z)=0\) locally for \(z\)?
\end{exercise}
\begin{hint}
Differentiate with respect to the variable being solved for.
\end{hint}
\begin{solution}
\(F_z\ne0\).
\end{solution}

\begin{exercise}[Regular level set]\label{exr:ii09-expand-05}
Is \(0\) a regular value of \(F(x,y)=x^2+y^2-1\)?
\end{exercise}
\begin{hint}
Check the gradient on the zero set.
\end{hint}
\begin{solution}
Yes; \(\nabla F=(2x,2y)\) never vanishes on the unit circle.
\end{solution}

\subsection*{Proofs}

\begin{exercise}[Inverse derivative formula]\label{exr:ii09-expand-06}
Derive \(D(f^{-1})(f(a))=Df(a)^{-1}\).
\end{exercise}
\begin{hint}
Differentiate \(f^{-1}\circ f=\mathrm{id}\).
\end{hint}
\begin{solution}
The chain rule gives \(D(f^{-1})(f(a))Df(a)=I\), so multiply by the inverse.
\end{solution}

\begin{exercise}[Implicit derivative formula]\label{exr:ii09-expand-07}
Derive \(Dg=-D_yF^{-1}D_xF\) from \(F(x,g(x))=0\).
\end{exercise}
\begin{hint}
Differentiate the identity and solve the resulting linear equation.
\end{hint}
\begin{solution}
\(D_xF+D_yF\,Dg=0\), hence the formula.
\end{solution}

\begin{exercise}[Local openness]\label{exr:ii09-expand-08}
Why does an invertible derivative imply \(f\) maps a sufficiently small neighborhood of \(a\) onto a neighborhood of \(f(a)\)?
\end{exercise}
\begin{hint}
Use the local diffeomorphism conclusion of the inverse theorem.
\end{hint}
\begin{solution}
The local inverse is defined on an open neighborhood of \(f(a)\), so that neighborhood lies in the image.
\end{solution}

\begin{exercise}[Regular graph dimension]\label{exr:ii09-expand-09}
If \(F:\mathbb R^{n+m}\to\mathbb R^m\) has full \(y\)-rank, why is the solution set locally \(n\)-dimensional?
\end{exercise}
\begin{hint}
The implicit theorem graphs \(y=g(x)\) over \(n\) free \(x\)-variables.
\end{hint}
\begin{solution}
A local parametrization \(x\mapsto(x,g(x))\) uses exactly \(n\) parameters.
\end{solution}

\subsection*{Counterexamples and hypothesis testing}

\begin{exercise}[Local not global]\label{exr:ii09-expand-10}
Give a smooth map with invertible derivative everywhere that is not globally injective.
\end{exercise}
\begin{hint}
Use exponential polar coordinates.
\end{hint}
\begin{solution}
\((x,y)\mapsto(e^x\cos y,e^x\sin y)\) has determinant \(e^{2x}>0\) but is \(2\pi\)-periodic in \(y\).
\end{solution}

\begin{exercise}[Inverse exists but theorem fails]\label{exr:ii09-expand-11}
Show \(f(x)=x^3\) has a global inverse although the inverse theorem fails at \(0\).
\end{exercise}
\begin{hint}
Compute \(f'(0)\) and the inverse.
\end{hint}
\begin{solution}
\(f'(0)=0\), so the theorem does not apply; nevertheless \(f^{-1}(y)=y^{1/3}\) exists, though its derivative blows up at \(0\).
\end{solution}

\begin{exercise}[Singular implicit block]\label{exr:ii09-expand-12}
For \(F(x,y)=y^2-x^2\), explain what fails at the origin when solving for \(y\).
\end{exercise}
\begin{hint}
Compute \(F_y(0,0)\) and inspect branches.
\end{hint}
\begin{solution}
\(F_y=2y=0\); the zero set has two branches \(y=\pm x\), so local uniqueness as one graph fails.
\end{solution}

\subsection*{Applications and investigations}

\begin{exercise}[Sensitivity of inverse]\label{exr:ii09-expand-13}
What does a large norm of \(Df(a)^{-1}\) suggest about local inversion?
\end{exercise}
\begin{hint}
The inverse derivative controls first-order amplification.
\end{hint}
\begin{solution}
Small output perturbations can require large input changes, indicating local ill-conditioning.
\end{solution}

\begin{exercise}[Newton correction]\label{exr:ii09-expand-14}
How is Newton's step related to the inverse theorem's linearization?
\end{exercise}
\begin{hint}
Solve the linearized equation \(Df(x_k)\Delta=-f(x_k)\).
\end{hint}
\begin{solution}
Newton uses the derivative inverse to correct the current residual, mirroring the theorem's local linear solvability mechanism.
\end{solution}

\subsection*{Challenge problems}

\begin{exercise}[Polar map local charts]\label{exr:ii09-expand-15}
For \(f(x,y)=(e^x\cos y,e^x\sin y)\), describe a neighborhood on which \(f\) is injective.
\end{exercise}
\begin{hint}
Restrict \(y\) to an interval shorter than \(2\pi\) and stay away from its endpoints.
\end{hint}
\begin{solution}
Any sufficiently small rectangle around a point, with \(y\)-width below \(2\pi\), gives a local chart; the inverse theorem guarantees such a neighborhood.
\end{solution}

\begin{exercise}[Implicit second derivative]\label{exr:ii09-expand-16}
For \(x^2+y^2=1\) near \((0,1)\), compute \(y''(0)\).
\end{exercise}
\begin{hint}
From \(y'=-x/y\), differentiate once more or use \(y=\sqrt{1-x^2}\).
\end{hint}
\begin{solution}
\(y''(0)=-1\).
\end{solution}
% END VOL02-EXPANSION II09-exercises-01
\section*{Chapter summary}
The chapter's tools now form part of the canonical Volume II analysis toolkit and will be used in later approximation, fixed-point, and convergence arguments.
